\chapter{总结与展望}

\section{全文工作总结}
本文围绕流程工业场景下的领域自适应故障诊断问题，面向 Tennessee Eastman Process Domain Adaptation 数据集，完成了从任务建模、统一实验脚手架构建、典型基线比较到改进方法设计与初步验证的一系列工作。全文首先对流程工业跨工况故障诊断问题进行了任务化表述，明确了单源到单目标无监督领域自适应的研究边界；随后构建了 manifest 驱动的数据组织方式、统一的 FCN 骨干网络、统一的优化配置和结构化结果产物；在此基础上，对 Source Only、CORAL、DAN、DANN、CDAN 与 DeepJDOT 等典型基线方法开展了系统比较；最后提出并实现了可靠性门控的类条件时序自适应方法 RCTA，并在当前已完成的代表性单源场景上完成了阶段性验证。

\section{主要结论}
本文的主要结论可概括为以下四点。

\begin{enumerate}
\item TEP 多工况故障诊断可以被清晰表述为单源到单目标的无监督领域自适应问题，并能够在统一数据划分、统一骨干网络和统一模型选择准则下开展可复现比较。
\item 在当前已完成的 4 个代表性单源场景上，领域自适应方法整体显著优于 Source Only，说明跨工况分布偏移确实需要显式对齐机制加以缓解。其中，在 6 种基线方法中，DANN 的平均目标域准确率最高，达到 0.8329；DAN 的平均表现为 0.8118，并展现出较好的跨场景稳定性。
\item 基线结果同时表明，单纯全局对齐并不能彻底解决多类别流程工业时序任务中的类别混叠与场景敏感问题。特别是在 Source Only 已经较强的场景中，若对齐压力控制不当，仍可能发生负迁移。
\item 本文提出的 RCTA 方法能够在保留强基线对齐能力的基础上，进一步利用可靠性门控伪标签、双原型记忆和一致性约束提升跨域同类聚合能力。当前已完成的四个代表性单源场景结果显示，RCTA 的平均目标域准确率达到 0.8590，相比 Source Only 平均提升 0.1884，相比最强基线平均再提升 0.0186。
\end{enumerate}

\section{不足与未来展望}
当前工作仍存在若干不足。首先，虽然当前 70 组实验计划中的 4 个单源场景已经全部完成，但 6 个多源场景仍在运行与整理过程中，因此当前正文中的系统结果仍主要围绕单源任务展开。其次，RCTA 已经完成初步验证，但严格意义上的模块消融、参数敏感性分析、统计显著性分析以及更系统的可视化解释仍有待补充。最后，本文目前聚焦于离线无监督领域自适应设置，对在线适应、持续诊断、多源异构融合和工程部署约束等问题尚未展开深入研究。

后续工作可以从以下几个方向继续推进。

\begin{enumerate}
\item 完成当前 70 组实验计划中多源部分的运行与汇总，形成覆盖单源与多源场景的完整结果表、热力图和统计分析；
\item 围绕 RCTA 开展严格消融实验，分别验证可靠性门控、双原型记忆、一致性正则、MCC 正则和基础对齐分支选择对最终性能的贡献；
\item 深入研究无标签模型选择问题，比较不同源域代理指标、目标熵指标、域判别器指标及其组合策略对最终 checkpoint 质量的影响；
\item 引入更多可解释性分析，如关键类别混淆矩阵对比、t-SNE 聚类结构变化、原型演化轨迹和目标伪标签质量统计，以增强方法分析深度；
\item 进一步探索面向工业应用的更复杂设置，例如多源到单目标迁移、在线适应、少样本目标域微调以及与工业先验知识结合的混合方法。
\end{enumerate}

总体来看，本文已经完成了一条较完整且可继续扩展的研究路线。随着后续多源实验、消融实验和附加分析的补充，本文有望从“形成统一基线与阶段性改进”进一步推进到“形成较完整的流程工业领域自适应故障诊断研究结果”。
